\documentclass{article}
\usepackage[T1]{fontenc}
\usepackage{lmodern}
\usepackage{graphicx}
\usepackage{amsmath,amssymb}
\usepackage[a4paper,margin=2cm]{geometry}
\usepackage[absolute,overlay]{textpos}
\setlength{\TPHorizModule}{1mm}
\setlength{\TPVertModule}{1mm}
\usepackage[indonesian]{babel}
\usepackage{hyperref}
\setlength{\emergencystretch}{2em}
% unit: Halaman 22

\begin{document}

% === Halaman 22 (llm) ===

\section*{Kategori \textit{cipher} berbasis bit}

\begin{enumerate}
    \item \textbf{Cipher Alir (\textit{Stream Cipher})}
    \begin{itemize}
        \item beroperasi pada bit (atau \textit{byte}) secara individual
        \item enkripsi/dekripsi pesan secara bit per bit (atau \textit{byte per byte}) hanya dengan operasi XOR saja
    \end{itemize}
    
    \item \textbf{Cipher Blok (\textit{Block Cipher})}
    \begin{itemize}
        \item beroperasi pada blok-blok bit (sekumpulan bit) \\ (contoh: 128-bit/blok = 16 byte/blok)
        \item enkripsi/dekripsi pesan secara blok per blok bit dengan komputasi yang lebih kompleks dan dilakukan berulang kali dalam sejumlah putaran.
    \end{itemize}
\end{enumerate}

\end{document}
